\SourceLeaf{150}
\begin{center}
\begin{tabular}{l@{\;}r}
in der Landdrostei Stade .\,.\,.\,.\,.\,.\,.\,.\,.\,.\,.\,. & 4,21\\
in der Landdrostei Aurich .\,.\,.\,.\,.\,.\,.\,.\,.\,.\,. & 3,28\\
in der Berghauptmannschaft Clausthal .\,.\,.\,. & 2,56
\end{tabular}
\end{center}
\noindent stehende Ehen gegen eine Witwe.

{\rightskip=0pt plus 1em
In der zweiten Hälfte des vorigen Jahrhunderts hat man manche Witwenkassen zu Grunde gehen sehen, weil sie auf die (obigen Bemerkungen zufolge ganz falsche) Voraussetzungen gegründet waren\footnote{Manche allerdings auch in Folge von \textit{noch} gröbern Fehlern.}, dass das Verhältniss der stehenden Ehen zu den Witwen, welches sich aus Zählungen für ein ganzes Land ergibt, auch für Witwenkassen als maassgebend betrachtet werden dürfe. \textsc{Kritter} war einer von denen, welche nicht ermüdeten, solche falsche Grundsätze zu bekämpfen. So oft er aber in seinen Schriften mit positiven eignen Angaben auftritt, heisst es \textit{entweder} nur, dass man \textit{zum allerwenigsten} Eine Witwe auf zwei Ehen rechnen müsse (z. B. \textsc{Kritter}’s Vorstellung des bisherigen Erfolgs u. s. w. S. 22); \textit{oder} er verwahrt sich ausdrücklich (Prufung eines Aufsatzes in der Berliner Monatsschrift S. 28), dass das Verhältniss \(2:1\) nur alsdann angenommen werden dürfe, wenn der sechste Theil aller Witwen sich wieder verheirathe, und ihre Pensionen damit erlöschen; \textit{oder} er gibt Zahlen an, die eben bedeutend \textit{mehr} Witwen als die halbe Zahl der stehenden Ehen erbringen; z. B.\par}

\begin{quote}
in der eben angeführten Schrift S. 27 ist eine Rechnung geführt, wonach gegen 3293 Ehepaare 1831 Witwen kommen,

Prüfung einer kleinen Schrift u. s. w. S. 17, wo auf 100 Ehemänner 60 Witwen gerechnet werden.
\end{quote}
\noindent Eben dieses Verhältniss ist aufgestellt in
\begin{quote}
Sammlung dreier Aufsätze über die Calenbergische, Preussische und Dänische Witwenversorgungsanstalten S. 30—35.

Sammlung wichtiger Erfahrungen u. s. w. S. 35, wo \textsc{Kritter} auf 450 Ehemänner 280 Witwen rechnet.
\end{quote}
\noindent\hspace*{\parindent}Uebrigens hat \textsc{Kritter} seine Angaben nirgends auf wirkliche directe Erfahrungen gestützt (dergleichen von genügender Art auch schwerlich aus Deutschland damals zu beschaffen waren), sondern auf die Mortalitätstafeln und auf gewisse Voraussetzungen rücksichtlich der durchschnittlichen Altersverhältnisse der in die Gesellschaft eintretenden, Voraussetzungen, die jedenfalls nicht ungünstiger gewählt sind, als sie sich bei unserer Professoren-Witwenkasse wirklich finden.

{\rightskip=0pt plus 1em
Wirkliche Erfahrungen aus einem ausgedehnten und mit unsrer Witwenkasse wohl ungefähr auf gleiche Linie zu stellenden Kreise finde ich in \textsc{Price} Observations on reversionary payments, S. 79. 269. 276 (nach der dritten Ausgabe dieses Werks von 1773), wo für die Gesammtheit der Pfarrer und Professoren in Schottland nach 17jährigem Durchschnitt die Zahl der stehenden Ehen zu 667, und die der Witwen zu 380 angegeben wird, also sehr nahe in dem Verhältniss von 7 zu 4. Es wird zugleich bemerkt, dass jene Standesklassen dort durchschnittlich mit dem Alter von 27 Jahren in den Genuss des Einkommens von ihren Stellen kamen, also gewiss früher, als durchschnittlich bei den Professoren deutscher Universitäten angenommen werden darf; auch waren bei jenen Witwen die Wiederverheirathungen nicht so selten, wie bei den Witwen Göttingischer Professoren.\par}


Aus den eignen directen Erfahrungen bei unsrer Gesellschaft lässt sich, auch abgesehen von der Kleinheit ihres Umfanges, schon deswegen das für sie gültige Normalverhältniss nicht ableiten, weil die Bewegung der Anzahl der stehenden Ehen in derselben eine unbekannte Grösse ist; diese Anzahl ist nemlich nur für zwei Zeitpunkte, aus der ganzen hundertjährigen Dauer der Anstalt, bekannt, für den gegenwärtigen Augenblick, und nach der P.’schen Angabe (oben S. 132) für den Zeitpunkt der Verhandlungen
